\documentclass[noauthor,noinstructornotes]{ximera}

\title{Week 2}
\author{Math 1193 Design Team}

\begin{document}
\begin{abstract}
    Activities for Week 2.
\end{abstract}
\maketitle

\begin{problem}
Say you sell shoes at \$120 per pair.

\begin{enumerate}
    \item How much do you 
    make if you sell 0 pairs of shoes? 1? 2? 3? 
    \item How much do you 
    make if you sell 32 pairs of shoes?
    \item How much do you make if you sell \(x\)
    pairs of shoes?
    \item Write the formula for a function \(R\)
    whih takes in a number of pairs of shoes and
    returns the revenue made.
    \item Graph your function \(R\).
    \item What does the slope of \(R\) represent
    in the problem?
\end{enumerate}

\begin{instructorNotes}
This is all about real-world examples of linear relationships,
emphasizing the fact that if students can answer the question 
for a concrete number, they can answer the question for 
variables. The last parts of this question is meant to 
get students used to function notation. 
\end{instructorNotes}

\end{problem}
Say it costs you \$30 in production costs to make a pair
of shoes and \$2000 per month to rent out the shoe-making
facility.

\begin{enumerate}
    \item How much does it cost to 
    make 0 pairs of shoes per month? 1? 2? 3? 
    \item How much does it cost to 
    make 12 pairs of shoes per month?
    \item How much does it cost to 
    make \(x\) pairs of shoes per month? 
    \item Write the formula for a cost function \(C\)
    whih takes in a number of pairs of shoes and
    returns the cost per month to make them.
    \item Graph your function \(C\).
    \item What does the slope of \(C\) represent
    in the problem?
\end{enumerate}

\begin{instructorNotes}
This is all about real-world examples of linear relationships,
emphasizing the fact that if students can answer the question 
for a concrete number, they can answer the question for 
variables. The last parts of this question is meant to 
get students used to function notation. 
\end{instructorNotes}
\begin{problem}

\end{problem}

\begin{problem}
Consider the function \(f(x) = \frac{1}{2}x + 5\).

\begin{enumerate}
    \item Write down a word problem about some
    quantity that can be represented by \(f(x)\).
    Be sure to define the inputs and outputs.
    \item What are the \(x\)- and \(y\)-intercepts
    of the function? What do those represent
    in your problem?
    \item What is the slope of \(f\)? What does 
    that represent in your problem?
\end{enumerate}

\begin{instructorNotes}
This is meant to draw on students' creativity. We'll see how that goes.
\end{instructorNotes}
\end{problem}

\begin{problem}
Say \(f(x) = 2x - g(x)\) and \(g(x) = 5 + x\).

\begin{enumerate}
    \item What is the slope of \(f\)?
    \item What is the \(y\)-intercept?
    \item What is \(f(2)\)?
    \item Graph \(f\).
\end{enumerate}

\begin{instructorNotes}
This is meant to gauge familiarity with substitution and function notation.
\end{instructorNotes}
\end{problem}

\end{document}
